\subsection{Eisenstein polynomials}

DEFINITION 2. --- Let A be a ring, p a prime ideal of A, and P a polynomial in A{[}T{]}. We say that P is an Eisenstein polynomial for p if it satisfies the following conditions:

\begin{enumerate}
\def\labelenumi{\alph{enumi})}
\item
  P is monic of degree \(d \geqslant 1\) ;
\item
  we have \(\mathrm{P(T)}\equiv \mathrm{T}^d\bmod \mathfrak{pA}[\mathrm{T}];\)
\item
  we have \(\mathbf{P}(0)\notin \mathfrak{p}^2\)
\end{enumerate}

In other words, an Eisenstein polynomial for p is a polynomial of the form \(\mathbf{P}(\mathbf{T}) = \mathbf{T}^{d} + \sum_{i=1}^{d} a_{i} \mathbf{T}^{d-i}\) , with \(d \geqslant 1\) , where \(a_{1}, \ldots, a_{d-1}\) belong to p and \(a_{d}\) lies in \(p - p^{2}\) .

We say that P is an Eisenstein polynomial for \(pA_{p}\) if the canonical image of P in the polynomial ring \(A_{p}[T]\) is an Eisenstein polynomial for the ideal \(pA_{p}\) ; this also means that P is an Eisenstein polynomial for p and that it satisfies in addition the following condition, stronger than c):

\(c^{\prime})\) every element \(a\) of A such that \(a\mathbf{P}(0)\in\mathfrak{p}^{2}\) belongs to \(\mathfrak{p}\) .

PROPOSITION 3. --- Let A be a ring, p a prime ideal of A, and P ∈ A{[}T{]} an Eisenstein polynomial for p.

\begin{enumerate}
\def\labelenumi{\alph{enumi})}
\item
  There is no decomposition of the form \(P = P_{1}P_{2}\) where \(P_{1}\) and \(P_{2}\) are two monic polynomials of A{[}T{]} distinct from 1.
\item
  Suppose A is integrally closed, with field of fractions K. Then P is irreducible in K[T].
\end{enumerate}

Let \(\varphi\) be the canonical homomorphism from A into the field of fractions \(k\) of A/p, and let \(\varphi':\mathrm{A}[\mathrm{T}]\to k[\mathrm{T}]\) be the extension of \(\varphi\) such that \(\varphi'(\mathrm{T})=\mathrm{T}\). Suppose we have \(\mathbf{P}=\mathbf{P}_{1}\mathbf{P}_{2}\) where \(\mathbf{P}_{1}\) and \(\mathbf{P}_{2}\) are two monic polynomials of A{[}T{]} distinct from 1. We then have \(\mathrm{T}^{d}=\varphi'(\mathbf{P}_{1})\varphi'(\mathbf{P}_{2})\) in \(k[\mathrm{T}]\) , denoting \(d\) the degree of P. If \(d_{i}\) is the degree of \(\mathbf{P}_{i}\) , one therefore has \(\varphi'(\mathbf{P}_{i})=\mathrm{T}^{d_{i}}\) , that is, \(\mathbf{P}_{i}(\mathrm{T})\equiv\mathrm{T}^{d_{i}}\text{ mod. pA[T]}\) , and in particular \(\mathbf{P}_{i}(0)\in\mathfrak{p}\) . But then \(\mathrm{P}(0)=\mathrm{P}_{1}(0).\mathrm{P}_{2}(0)\) belongs to \(\mathfrak{p}^{2}\) contrary to the hypotheses. This proves \(a\).

The assertion \(b)\) follows from \(a)\) and from prop. 11 of V, § 1, no. 3.

Let A be a Noetherian local ring and \(P_{1}, \ldots, P_{r}\) monic polynomials in A{[}T{]}, of degree \(\geqslant 2\). Let q be the ideal of A{[}T \(_{1}\) , \ldots, T \(_{r}\) {]} generated by \(\mathrm{P}_{1}(\mathrm{T}_{1}), \ldots, \mathrm{P}_{r}(\mathrm{T}_{r})\) and let B be the quotient A-algebra A{[}T \(_{1}\) , \ldots, T \(_{r}\) {]}/q. For \(1 \leqslant i \leqslant r\) , denote by \(d_{i}\) the degree of \(P_{i}, t_{i}\) the class of T \(_{i}\) modulo q, and \(\gamma_{i}\) the class of \(c_{i} = \mathrm{P}_{i}(0)\) modulo \(m_{A}^{2}\). We assume that we have \(\mathrm{P}_{i}(\mathrm{T}) \equiv \mathrm{T}^{d_{i}} \mod \mathrm{m}_{\mathrm{A}}\mathrm{A}[\mathrm{T}]\) for \(1 \leqslant i \leqslant r\) .

PROPOSITION 4. --- a) The ring B is local and Noetherian, with maximal ideal

\[
\mathfrak {m} _ {\mathrm{B}} = \mathbf {B m} _ {\mathrm{A}} + \sum_ {i = 1} ^ {r} \mathbf {B t} _ {i}.
\]

We have \(\dim(\mathbf{A}) = \dim(\mathbf{B})\) and \([\kappa_{\mathbf{B}}: \kappa_{\mathbf{A}}] = 1\). The monomials \(t_1^{\alpha(1)} \ldots t_r^{\alpha(r)}\) , with \(0 \leqslant \alpha(i) < d_i\) for \(1 \leqslant i \leqslant r\), form a basis of the A-module B.

\begin{enumerate}
\def\labelenumi{\alph{enumi})}
\setcounter{enumi}{1}
\item
  Let \(\lambda\) be the canonical homomorphism from \(\mathfrak{m}_{\mathbf{A}} / \mathfrak{m}_{\mathbf{A}}^{2}\) into \(\mathfrak{m}_{\mathbf{B}} / \mathfrak{m}_{\mathbf{B}}^{2}\). Then the kernel of \(\lambda\) is the \(\kappa_{\mathbf{A}}\) -vector space generated by \(\gamma_{1}, \ldots, \gamma_{r}\) . The classes of the elements \(t_{1}, \ldots, t_{r}\) form a basis over \(\kappa_{\mathbf{A}}\) of the cokernel of \(\lambda\) .
\item
  For B to be regular, it is necessary and sufficient that A be regular and that \(\gamma_{1},\ldots,\gamma_{r}\) are linearly independent in the \(\kappa_{\mathbf{A}}\) -vector space \(\mathfrak{m}_{\mathbf{A}}/\mathfrak{m}_{\mathbf{A}}^{2}\) .
\end{enumerate}

The A-algebra B is isomorphic to the tensor product \(B_{1} \otimes_{A} \cdots \otimes_{A} B_{r}\) with \(B_{i} = A[T]/(P_{i})\) for \(1 \leqslant i \leqslant r\). It follows that the monomials \(t_{1}^{\alpha(1)} \ldots t_{r}^{\alpha(r)}\) , with \(0 \leqslant \alpha(i) < d_{i}\) for \(1 \leqslant i \leqslant r\) , form a basis of the A-module B. In particular, B is integral over A, so A and B have the same dimension by th. 1 of § 2, no. 3.

By cor. 3 of prop. 9 of IV, § 2, no. 5, the ring B is Noetherian, and every maximal ideal of B contains B.m\_A. Moreover, in view of the hypothesis made on P\_1, \ldots P\_r and the relation P\_i(t\_i) = 0, we have t\_i\^{}\{d\_i\} ∈ B.m\_A for 1 ≤ i ≤ r. Hence every maximal ideal of B contains t\_1, \ldots, t\_r, hence also the ideal q' = B.m\_A + Bt\_1 + \ldots{} + Bt\_r. Now we have m\_A = A ∩ q' and B = A + q', so B/q' is isomorphic to A/m\_A, and q' is a maximal ideal of B; consequently, B is local and we have {[}κ\_B : κ\_A{]} = 1. This proves a).

Set \(\mathfrak{r} = \mathfrak{m}_{A}^{2} + \sum_{i=1}^{r} A c_i\) , and let \(\varphi\) be the canonical homomorphism from \((A/\mathfrak{m}_{A}^{2})[T_1, ..., T_r]\) onto \(B/\mathfrak{m}_{B}^{2}\) . Since one has \(m_B = B.m_A + \sum_{i=1}^{r} B t_i\) , the kernel \(n\) of \(\varphi\) is the ideal generated by the classes \(\overline{P}_i(T_i)\) of the polynomials \(P_i(T_i)\) modulo \(\mathfrak{m}_{A}^{2}.A[T_1, ..., T_r]\) and the monomials \(T_i T_j\) and \(x T_i\) for \(1 \leqslant i, j \leqslant r\) and \(x\) in \(m_A/m_A^2\) . By the hypothesis made on \(P_i\) , namely \(P_i(T) \equiv T^{d_i} \mod. m_A.A[T]\) , one may replace \(\overline{P}_i(T_i)\) by \(\gamma_i\) in this description of \(n\) ; consequently, the ring \(B/m_B^2\) is isomorphic to the quotient of \(A/\mathfrak{r}\) \([T_{1}, ..., T_{r}]\) by the graded ideal generated by the monomials \(T_{i}T_{j}\) and \(xT_{i}\) for \(x\) in \(m_A/\mathfrak{r}\) . Denote by \(\tau_{i}\) the class of \(t_{i}\) modulo \(m_{B}^{2}\) ; we thus have

\[
\mathrm{B} / \mathfrak {m} _ {\mathrm{B}} ^ {2} = (\mathrm{A} / \mathfrak {r}) \oplus \kappa_ {\mathrm{A}} \tau_ {1} \oplus \dots \oplus \kappa_ {\mathrm{A}} \tau_ {r},\tag{5}
\]

whence

\[
\mathfrak {m} _ {\mathrm{B}} / \mathfrak {m} _ {\mathrm{B}} ^ {2} = (\mathfrak {m} _ {\mathrm{A}} / \mathfrak {r}) \oplus \kappa_ {\mathrm{A}} \tau_ {1} \oplus \dots \oplus \kappa_ {\mathrm{A}} \tau_ {r}.\tag{6}
\]

Assertion b) follows immediately from that.

By formula (6) and the relation \([\kappa_{\mathrm{B}}:\kappa_{\mathrm{A}}] = 1\) , one has

\[
\left[ \mathfrak {m} _ {\mathrm{B}} / \mathfrak {m} _ {\mathrm{B}} ^ {2}: \kappa_ {\mathrm{B}} \right] = \left[ \mathfrak {m} _ {\mathrm{A}} / \mathfrak {m} _ {\mathrm{A}} ^ {2}: \kappa_ {\mathrm{A}} \right] + \left\{r - \left[ \mathfrak {r} / \mathfrak {m} _ {\mathrm{A}} ^ {2}: \kappa_ {\mathrm{A}} \right] \right\}.\tag{7}
\]

Now, the \(\kappa_{\mathrm{A}}\) -vector space \(\mathfrak{r}/\mathfrak{m}_{\mathrm{A}}^{2}\) is generated by \(\gamma_{1}, \ldots, \gamma_{r}\) and one has

\[
\dim (\mathbf {B}) = \dim (\mathbf {A}) \leqslant [ \mathfrak {m} _ {\mathbf {A}} / \mathfrak {m} _ {\mathbf {A}} ^ {2}: \kappa_ {\mathbf {A}} ].\tag{8}
\]

Assertion (c) follows immediately from formulas (7) and (8).

COROLLARY. --- Let A be a regular Noetherian local ring and let P ∈ A{[}T{]} an Eisenstein polynomial for m\_A. The ring B = A{[}T{]}/(P) is a regular local Noetherian ring, of the same dimension as A, and one has {[}κ\_A : κ\_B{]} = 1. Finally, we have m\_B = B.m\_A + Bt, where t is the class of T modulo (P).

The case where P has degree ≥ 2 follows from Prop. 4, taking r = 1; when P has degree 1, that is, of the form T - c with \(c \in m_{A}\); the corollary is immediate.

PROPOSITION 5. --- Let A be an integral domain, with field of fractions K, and let L be a finite algebraic extension of K. We denote by B the integral closure of A in L, and let p be a prime ideal of A.

Suppose that the local ring \(\mathbf{A}_{\mathfrak{p}}\) is Noetherian and regular; let \(t\) be an element of \(\mathbf{L}\) such that \(\mathbf{L} = \mathbf{K}(t)\), and suppose that there exists in \(\mathbf{A}[\mathbf{T}]\) an element \(\mathbf{P}\), an Eisenstein polynomial for \(\mathfrak{p}\mathbf{A}_{\mathfrak{p}}\), of which \(t\) is a root.

\begin{enumerate}
\def\labelenumi{\alph{enumi})}
\item
  There exists in B a unique prime ideal q lying over p.
\item
  The local ring \(\mathbf{B}_{\mathfrak{q}}\) is Noetherian and regular, of the same dimension as \(\mathbf{A}_{\mathfrak{p}}\) .
\item
  We have \(\mathbf{B}_{\mathfrak{q}} = \mathrm{A}_{\mathfrak{p}}[t]\) .
\item
  The canonical homomorphism from A/p to B/q induces an isomorphism of the fields of fractions of these rings.
\end{enumerate}

Set \(C = A_{p}[t]\) and denote by \(d\) the degree of P. By prop. 3 applied to the ring \(A_{p}\) the Eisenstein polynomial P is irreducible in K{[}T{]} and \((1, t, ..., t^{d-1})\) is a basis of L over K, hence of C over \(A_{p}\) . Since P is monic, the kernel of the canonical homomorphism from \(A_{p}[T]\) onto C is equal to (P). By the corollary of prop. 4 above, C is therefore a regular noetherian local ring of the same dimension as \(A_{p}\) , the maximal ideal \(m_{C}\) of C is generated by \(p \cup \{t\}\) and the field \(\kappa_{C}\) is a trivial extension of the field of fractions of A/p. To prove prop. 5, it therefore suffices to show that there exists a unique prime ideal q of B above p, and that one has \(C = B_{q}\) .

Set \(S = A - p\) . We know (V, § 1, no. 5, prop. 16) that the integral closure of

\(A_{p}\) in L is equal to \(S^{-1}B\) . Moreover, t is integral over \(A_{p}\) , and the ring \(C = A_{p}[t]\) is a regular noetherian local ring, hence integrally closed (no. 2, cor. 1 of th. 1). We therefore have \(C = S^{-1}B\) . Consequently, the ring \(S^{-1}B\) is local and has a unique maximal ideal. By V, § 2, no. 1, prop. 1, there exists a unique prime ideal of \(S^{-1}B\) above \(pA_{p}\) , so B has a unique prime ideal q above p (loc. cit., lemma 1), and one has \(B_{q} = S^{-1}B = C\) .

COROLLARY. --- Suppose that \(A_{p}\) is a discrete valuation ring. Then \(B_{q}\) is a discrete valuation ring, t is a uniformizer of \(B_{q}\) , and one has

\[
f (\mathrm{B} _ {\mathfrak {q}} / \mathrm{A} _ {\mathfrak {p}}) = 1, e (\mathrm{B} _ {\mathfrak {q}} / \mathrm{A} _ {\mathfrak {p}}) = [ \mathrm{L}: \mathrm{K} ]\tag{9}
\]

(VI, § 8, no. 1).

Indeed, discrete valuation rings are precisely the regular noetherian local rings of dimension 1; setting \(d = [\mathrm{L}:\mathrm{K}]\) , one has \(t^{d}\in\mathfrak{m}_{\mathrm{A}}-\mathfrak{m}_{\mathrm{A}}^{2}\) , whence \(d=e(\mathrm{B}_{\mathfrak{q}}/\mathrm{A}_{\mathfrak{p}})\) . We have \([\kappa_{\mathrm{B}}:\kappa_{\mathrm{A}}]=1\) , whence \(f(\mathrm{B}_{\mathfrak{q}}/\mathrm{A}_{\mathfrak{p}})=1\) .

Examples. --- 1) Set A = Z and L = Q( \(p^{1/d}\) ), where p is a prime number and d an integer ≥ 2. Let B denote the integral closure of Z in L. Since the polynomial T \(^{d}\) - p in Z{[}T{]} is an Eisenstein polynomial for pZ \(_{(p)}\) , there exists a unique prime ideal q of B above pZ. There therefore exists a unique normalized discrete valuation v of the field Q( \(p^{1/d}\) ) such that v(p) \textgreater{} 0; we have {[}L:K{]} = v(p) = d, and B/q is a field with p elements. The ring B\(_{q}\) of the valuation v is equal to Z \(_{(p)}\) {[} \(p^{1/d}\) {]}.

\begin{enumerate}
\def\labelenumi{\arabic{enumi})}
\setcounter{enumi}{1}
\tightlist
\item
  Set A = Z and \(L = R_{p^{f}}(\mathbf{Q})\) where p is a prime number and f an integer \(\geqslant 1\) (cf. A, V, p. 78). We therefore have \(L = \mathbf{Q}(\zeta)\) with \(\zeta = \exp(2\pi i/p^{f})\) . Let B be the integral closure of Z in L and P the polynomial of Z{[}T{]} such that
\end{enumerate}

\[
\mathrm{P} (\mathrm{T} - 1) = \left(\mathrm{T} ^ {p^{f}} - 1\right) / \left(\mathrm{T} ^ {p^{f - 1}} - 1\right).
\]

Set \(d = p^{f} - p^{f - 1}\) . We have \(\mathrm{P}(\zeta - 1) = 0\) , \(\mathrm{P}(0) = p\) and

\[
\mathrm{P} (\mathrm{T} - 1) \equiv (\mathrm{T} - 1) ^ {d} \bmod . p \mathbf {Z} [ \mathrm{T} ],
\]

whence P(T) ≡ T \(^{d}\) . Consequently, P is an Eisenstein polynomial for pZ \(_{(p)}\) . There is therefore a unique prime ideal q of B above pZ, and we have B \(_{q}\) = Z \(_{(p)}\) {[}ζ{]}; moreover, ζ - 1 is a uniformizer of B \(_{q}\) and one has

\[
[ \mathrm{L}: \mathrm{K} ] = d = p ^ {f} - p ^ {f - 1}.
\]

If \(v\) is the unique normalized valuation of \(\mathbf{Q}(\zeta)\) such that \(v(p)>0\) , one has \(v(p)=d\) . Moreover, the field \(\mathbf{B}/\mathfrak{q}\) has \(p\) elements. One can prove (cf. p. 96, exerc. 13) that \(\mathbf{B}\) is equal to \(\mathbf{Z}[\zeta]\) .



